\documentclass[a4paper,12pt]{article}
\usepackage[utf8]{inputenc}
\usepackage[ngerman]{babel}
\usepackage{amsmath}
\usepackage{parskip}

\setlength{\parindent}{0pt}

\title{Modul 294 - Aufgabenblatt 10}
\author{von Lukas Meier}
\date{Unterricht vom 03.03.2026}

\begin{document}

\maketitle

\section{Ziel}

Diese Lektion dient als Repetition der bisherigen Inhalte zur Vorbereitung auf die Prüfung in der folgenden Woche.

Nutzen Sie Ihre bisherigen Unterlagen, die vorhandenen Beispielprojekte und Ihren aktuellen Wissensstand aus den Lektionen 1 bis 9.

\section{Teil 1: Theorie kurz zusammenfassen}

Erstellen Sie eine eigene Zusammenfassung zu folgenden Themen:
\begin{itemize}
    \item Variablen, Datentypen, Arrays, Objekte
    \item Bedingungen, Schleifen, Funktionen
    \item DOM-Zugriff und Event-Handling
    \item Klassen, Module und saubere Struktur
    \item \texttt{fetch}, \texttt{async}/\texttt{await}, Fehlerbehandlung
    \item CRUD-Ablauf im Projekt
\end{itemize}

\section{Teil 2: Begriffe sicher beherrschen}

Erklären Sie zu jedem Begriff in 1 bis 2 Sätzen, wozu er dient:
\begin{itemize}
    \item \texttt{addEventListener}
    \item \texttt{querySelector}
    \item \texttt{await}
    \item \texttt{response.ok}
    \item JSON
    \item CRUD
\end{itemize}

\section{Teil 3: Code lesen}

Betrachten Sie den folgenden Code und beantworten Sie:

\begin{verbatim}
function renderNotes(notes) {
    notesList.innerHTML = "";

    for (const note of notes) {
        const li = document.createElement("li");
        li.textContent = `${note.title}: ${note.content}`;
        notesList.append(li);
    }
}
\end{verbatim}

\begin{itemize}
    \item Welche Daten kommen rein?
    \item Was macht die Funktion Schritt für Schritt?
    \item Welche DOM-Elemente werden verändert/erstellt?
    \item Wo könnte ein typischer Fehler entstehen?
\end{itemize}

\section{Teil 4: Mini-Aufgabe DOM}

Erstellen Sie ein kleines Formular mit einem Eingabefeld und einem Button.

Beim Klick auf den Button soll:
\begin{itemize}
    \item der Wert aus dem Input gelesen werden
    \item geprüft werden, ob er nicht leer ist
    \item ein neuer Listeneintrag im DOM erzeugt werden
\end{itemize}

\section{Teil 5: Mini-Aufgabe Async}

Skizzieren oder implementieren Sie eine Funktion \texttt{loadData()}, die:
\begin{itemize}
    \item Daten mit \texttt{fetch(...)} lädt
    \item den HTTP-Status prüft
    \item das Resultat als JSON verarbeitet
    \item Fehler sauber behandelt
\end{itemize}

\section{Teil 6: Mini-Aufgabe CRUD}

Beschreiben Sie für Ihre Notizen-App den Ablauf für:
\begin{itemize}
    \item Create
    \item Read
    \item Update
    \item Delete
\end{itemize}

Notieren Sie jeweils:
\begin{itemize}
    \item welcher Request ausgeführt wird
    \item welche Daten gesendet werden
    \item wie das UI danach aktualisiert wird
\end{itemize}

\section{Teil 7: Prüfungssimulation}

Bearbeiten Sie ohne Hilfe eine kurze Übung:
\begin{itemize}
    \item Event auf ein Formular binden
    \item Eingaben validieren
    \item eine asynchrone Funktion aufrufen
    \item Ergebnis im DOM anzeigen
\end{itemize}

Arbeiten Sie dabei bewusst in kleinen, nachvollziehbaren Schritten.

\section{Teil 8: Persönliche Checkliste}

Notieren Sie drei Themen, die Sie bereits sicher beherrschen, und drei Themen, die Sie vor der Prüfung nochmals anschauen wollen.

\section{Abschluss}

Die Lektion ist abgeschlossen, wenn:
\begin{itemize}
    \item Sie alle Themenfelder einmal wiederholt haben
    \item Sie mindestens eine kleine DOM- und eine kleine Async-Aufgabe gelöst haben
    \item Sie eine persönliche Prüfungs-Checkliste erstellt haben
\end{itemize}

\end{document}
